% Part: many-valued-logic
% Chapter: sequent-calculus
% Section: introduction

\documentclass[../../../include/open-logic-section]{subfiles}

\begin{document}

\olfileid{mvl}{seq}{int}

\olsection{引言}

经典逻辑的相继式演算是一种高效且简单的推导系统。
如果一个多值逻辑由有限多个真值的矩阵定义，即 $V$ 是有限的，
则可以为它提供一种相继式演算。
其思路来源于对经典情况下相继式的含义及推理规则形式的考察。

回想一下，一个相继式
\begin{align*}
    !A_1, \dots, !A_n & \Sequent !B_1, \dots, !B_n
\intertext{可以解释为公式}
    (!A_1 \land \cdots \land !A_m) & \lif (!B_1 \lor \cdots \lor
!B_n)
\end{align*}
换言之，赋值~$\pAssign{v}$ \emph{满足}相继式 $\Gamma \Sequent \Delta$，
当且仅当存在 $!A \in \Gamma$ 使得 $\pValue{v}(!A) = \False$，
或者存在 $!A \in \Delta$ 使得 $\pValue{v}(!A) = \True$。
根据这一解释，初始相继式 $!A \Sequent !A$ 总是被满足，因为或者 $\pValue v(!A) = \True$，
或者 $\pValue(!A) = \False$。

以下是 $\Log{LK}$ 中条件句的推理规则，其中省略了侧公式 $\Gamma$、$\Delta$：

\begin{defish}
    \Axiom$ \fCenter !A$
    \Axiom$ !B \fCenter $
    \RightLabel{\LeftR{\lif}}
    \BinaryInf$ !A \lif !B \fCenter $
    \DisplayProof
    \hfill
    \Axiom$ !A \fCenter !B$
    \RightLabel{\RightR{\lif}}
    \UnaryInf$ \fCenter !A \lif !B $
    \DisplayProof
\end{defish}

如果我们运用上述相继式的语义解释，就可以将 $\LeftR{\lif}$ 规则读作：
如果 $\pValue v(!A) = \True$ 且 $\pValue v(!B) = \False$，那么 $\pValue v(!A \lif !B) = \False$。
类似地，$\RightR{\lif}$ 规则说的是：如果 $\pValue v(!A) = \False$
或者 $\pValue v(!B) = \True$，那么 $\pValue v(!A \lif !B) =
\True$。
事实上，这些条件句其实是双条件句。
在标准形式下，$\LeftR{\land}$ 和 $\RightR{\lor}$ 规则对应的条件句不会是双条件句。
但存在这些规则的替代版本，在其中它们是双条件句：

\begin{defish}
    \Axiom$!A, !B, \Gamma \fCenter \Delta$
    \RightLabel{\LeftR{\land}}
    \UnaryInf$!A \land !B, \Gamma \fCenter \Delta$
    \DisplayProof
    \hfill
    \Axiom$ \Gamma \fCenter \Delta, !A, !B$
    \RightLabel{\RightR{\lor}}
    \UnaryInf$ \Gamma \fCenter \Delta, !A \lor !B$
    \DisplayProof
\end{defish}

将这一基本思想应用于 $n$ 值逻辑，便会得到一个具有 $n$ 个（而非两个）位置的相继式演算，
每个位置对应一个真值。对于具有 $V = \{\False, \Undef, \True\}$ 的三值逻辑，
相继式是表达式 $\Gamma \mid \Pi \mid \Delta$。
它在赋值~$\pAssign v$ 下被满足，当且仅当存在 $!A \in \Gamma$ 使得 $\pValue{v}(!A) = \False$，
或者存在 $!A \in \Delta$ 使得 $\pValue{v}(!A) = \True$，
或者存在 $!A \in \Pi$ 使得 $\pValue{v}(!A) = \Undef$。
因此，初始相继式 $!A \mid !A \mid !A$ 总是被满足。

\end{document}